311 records measure the propensity to report as much as the incidence of
problems: communities differ systematically in access, trust, language,
and expectation that reporting leads to action
\citep{kontokosta2017equity, kontokosta2021bias}, and prediction-driven
prioritization can institutionalize whatever bias the data carry
\citep{samorani2022overbooked}. Three design consequences follow in this
study. First, the target variable is named throughout as \emph{reported}
demand, and no result is interpreted as a measurement of need. Second,
because our spatial unit is the whole city, the most direct channel for
reporting-bias amplification --- reallocating resources \emph{away} from
low-reporting neighbourhoods --- is outside the action space of the
simulated planner; the family-level channel remains, and we report
per-family served fractions and unserved-stock shares for every policy
(supplement) so that systematic deprioritization of any family is visible
rather than hidden in aggregate scores. These tradeoffs are real: under
the equal request-equivalent objective the forecast-driven
full-distribution policy attains low aggregate loss while serving a
near-zero fraction of the low-volume water/sewer family in all four cities
(and of public safety in San Francisco), because minimizing total
request-equivalent stock concentrates capacity on high-volume families.
Lower aggregate simulated loss is therefore not a fairness guarantee.
Third, the priority weights $\omega_s$ encode a
value judgment, not a fact; we expose them as configuration, vary them over a pre-specified grid
in sensitivity analysis, and regard their selection in any real
deployment as a decision that belongs to accountable public officials,
not to a model. A forecast-driven allocation system used in practice
would additionally require visibility into where reporting data are
likely incomplete --- a measurement problem this paper's data cannot
solve and that we therefore flag as a limitation rather than simulate.
